%ID BT_3_1_3_01
\textbf{Bước 1. Miền hỗ trợ.} Nếu $|x|>M+at$, hai cận $x\pm at$ và toàn đoạn giữa chúng cùng nằm ngoài $[-M,M]$. D'Alembert cho $u=0$; các đạo hàm cũng bằng không ngoài miền ấy khi dữ liệu đủ trơn.

\textbf{Bước 2. Đạo hàm năng lượng.} Đặt $I(t)=\int_{\mathbb R}(a^2u_x^2+u_t^2)dx$. Trên mỗi khoảng thời gian hữu hạn có thể tích phân trên một đoạn cố định đủ lớn, vì sóng có hỗ trợ hữu hạn. Dùng PDE:
\[I'=2a^2\int(u_xu_{xt}+u_tu_{xx})dx=2a^2[u_tu_x]_{-\infty}^{+\infty}=0.\]
\textbf{Bước 3. Giá trị đầu.} Tại $t=0$, $u_x=f'$, $u_t=g$, nên
\[I(t)=I(0)=\int_{\mathbb R}(a^2(f')^2+g^2)dx.\]
Điều kiện hỗ trợ bảo đảm hạng biên bằng không, không phải chỉ vì ký hiệu cận là vô hạn. Với dữ liệu năng lượng, kết quả tiếp tục đúng bằng xấp xỉ dữ liệu trơn.

\textbf{Chi tiết tích phân từng phần trong đạo hàm năng lượng.} Đặt $E=\frac12\int_{-R}^R(u_t^2+a^2u_x^2+\kappa u^2)dx$. Với $u_{tt}-a^2u_{xx}+\gamma u_t+\kappa u=F$,
\[\begin{aligned}
E'&=\int_{-R}^R(u_tu_{tt}+a^2u_xu_{xt}+\kappa uu_t)dx\\
&=\int_{-R}^Ru_t(a^2u_{xx}-\gamma u_t-\kappa u+F)dx
+a^2\int_{-R}^Ru_xu_{xt}dx+\kappa\int_{-R}^Ruu_tdx\\
&=a^2[u_tu_x]_{-R}^R-a^2\int_{-R}^Ru_{tx}u_xdx
+a^2\int_{-R}^Ru_xu_{xt}dx-\gamma\int_{-R}^Ru_t^2dx+\int_{-R}^RFu_tdx\\
&=a^2[u_tu_x]_{-R}^R-\gamma\int_{-R}^Ru_t^2dx+\int_{-R}^RFu_tdx.
\end{aligned}\]
Hai hạng $\kappa uu_t$ triệt tiêu ngay khi thế PDE; hai tích phân $u_xu_{xt}$ triệt tiêu sau tích phân từng phần. Chọn $R$ lớn hơn miền hỗ trợ trên khoảng thời gian xét, nên $u_t=u_x=0$ ở $x=\pm R$. Với $I=2E$, nhân đẳng thức đạo hàm với $2$ để có $I'=0$.

%ID BT_3_1_3_02
\textbf{Bước 1. Tính sóng bằng nguyên hàm.} Dùng $H'=g$ và half-sum của $f$. Công thức tổng quát trong phần kiến thức cho cách kiểm tra $u_t(x,0)=g(x)$; dưới đây ghi rõ nguyên hàm và kết quả từng câu.

(i) $a=3$, $H=-\cos2s/2$. Hiệu hai cos bằng $2\sin2x\sin6t$, nên
\[u=\cos x\cos3t+\tfrac16\sin2x\sin6t,\qquad T=2\pi/3.\]
(ii) $a=5$, $H=s^2/2$. Hiệu bình phương bằng $20xt$, cho $u=\sin3x\cos15t+xt$, không tuần hoàn.

(iii) $a=2$, $H=\sin s$. Half-sum của bình phương bằng $x^2+4t^2$, nên $u=x^2+4t^2+\cos x\sin2t/2$, không tuần hoàn.

(iv) $a=6$, $H=s/2+\sin2s/4$ vì $\cos^2s=(1+\cos2s)/2$:
\[u=\tfrac12(e^{-|x-6t|}+e^{-|x+6t|})+t/2+\cos2x\sin12t/24.\]
Hạng $t/2$ ngăn tuần hoàn; li độ có góc tại $0$ nên không $C^1$ trên $x=\pm6t$.

(v) $a=1$, $g=(3\sin s-\sin3s)/4$, $H=-3\cos s/4+\cos3s/12$. Suy
\[u=\sin x\cos t-\tfrac12-\tfrac12\cos2x\cos2t+\tfrac34\sin x\sin t-\tfrac1{12}\sin3x\sin3t.\]
Mode tần số $1$ khác không, nên chu kỳ cơ bản là $2\pi$.

(vi) $a=7$, $f=1$. Nguyên hàm của xung tam giác được nối liên tục:
\[H(s)=\begin{cases}0&s\le-1,\\(s+1)^2/2&-1<s\le0,\\1-(1-s)^2/2&0<s<1,\\1&s\ge1.\end{cases}\quad
u=1+\frac{H(x+7t)-H(x-7t)}{14}.\]
Tại mỗi $x$ cố định, sau khi cả hai cận ra ngoài hỗ trợ, $u=15/14$, khác giá trị đầu $1$; vì thế không tuần hoàn. Nghiệm $C^1$, các đạo hàm cấp hai gãy trên $x\pm7t=-1,0,1$.

\textbf{Bước 2. Dựng đồ thị.} Sóng phải $A=f/2-H/(2a)$ được dịch sang phải $at$, sóng trái $B=f/2+H/(2a)$ sang trái $at$. Các hình kèm theo gồm lát cắt thời gian; những vị trí gãy được suy từ các cận chạm điểm gãy của dữ liệu. Với các câu trơn, đạo hàm từng mode xác nhận PDE; với câu (iv), (vi), chỉ khẳng định cổ điển ngoài các đường gãy đã chỉ ra.

\textbf{Chi tiết phép thế nguyên hàm và lượng giác.} Ở (i),
\[\frac{H(x+3t)-H(x-3t)}6
=\frac{\cos(2x-6t)-\cos(2x+6t)}{12}
=\frac{\sin2x\sin6t}{6}.\]
Ở (ii), $[(x+5t)^2-(x-5t)^2]/20=xt$. Ở (iii), $[\sin(x+2t)-\sin(x-2t)]/4=\cos x\sin2t/2$. Ở (iv), phần tuyến tính của $H$ cho $(12t/2)/12=t/2$; phần $\sin2s/4$ cho $2\cos2x\sin12t/48=\cos2x\sin12t/24$.

Ở (v), $f(s)=\sin s-\cos^2s=\sin s-1/2-\cos2s/2$, nên half-sum là $\sin x\cos t-1/2-\cos2x\cos2t/2$. Từ $\sin^3s=(3\sin s-\sin3s)/4$, nguyên hàm đã ghi cho
\[\frac{H(x+t)-H(x-t)}2=\frac34\sin x\sin t-\frac1{12}\sin3x\sin3t.\]
Ở (vi), trên $-1<s<0$, $H=\int_{-1}^s(r+1)dr=(s+1)^2/2$; trên $0<s<1$, $H=H(0)+\int_0^s(1-r)dr=1/2+s-s^2/2=1-(1-s)^2/2$. Sau $1$ giữ $H(1)=1$, bảo đảm không có bước nhảy của nguyên hàm.

%ID BT_3_1_3_03
\textbf{Bước 1. Tách sóng li độ và vận tốc.} $a=5$, $f=10\chi_{[-1,1]}$, $g=\chi_{[-1,1]}$. Tại $x=15$, cận phải $15+5t$ luôn ngoài hỗ trợ. Cận trái $15-5t$ nằm trong đoạn khi $14/5\le t\le16/5$. Đặt $\ell(t)=|[15-5t,15+5t]\cap[-1,1]|$.

\textbf{Bước 2. Viết áp suất theo từng giai đoạn.}
\[u(15,t)=\begin{cases}0&0\le t<14/5,\\5+(5t-14)/10&14/5\le t\le16/5,\\1/5&t>16/5.\end{cases}\]
Hạng $5$ là một nửa của li độ $10$; hạng thứ hai là $\ell/(2a)$. Ở hai đầu mút, giá trị phụ thuộc quy ước chỉ thị đóng; các vết một phía vẫn như trên.

\textbf{Bước 3. So sánh ngưỡng.} Trong khoảng giữa, hàm tăng, lớn nhất $5+2/10=26/5=5.2$, thấp hơn ngưỡng $8$. Sau xung chỉ còn $1/5$. Do đó áp suất không vượt ngưỡng. Không thể chỉ xét tích phân vận tốc mà bỏ sóng li độ gây bước nhảy $5$.

%ID BT_3_1_3_04
\textbf{Bước 1. Xung của nguồn tại thời điểm $T$.} Gọi $w(x,\tau;T)$ là nghiệm thuần nhất với $w(x,0;T)=0$, $w_\tau(x,0;T)=F(x,T)$. D'Alembert cho
\[w(x,\tau;T)=\frac1{2a}\int_{x-a\tau}^{x+a\tau}F(y,T)dy.\]
\textbf{Bước 2. Cộng các xung.} Với $V=\int_0^t w(x,t-T;T)dT$, Leibniz cho $V_{tt}=F(x,t)+\int_0^t w_{\tau\tau}dT$, còn $a^2V_{xx}=\int_0^ta^2w_{xx}dT$. Suy $V_{tt}-a^2V_{xx}=F$ và dữ liệu đầu bằng không.

\textbf{Bước 3. Khôi phục dữ liệu đầu.} Nếu đề có $f,g\ne0$ thì $\int w$ chỉ là phần nguồn; phải thêm $u_0$ của d'Alembert. Miền tích phân là tam giác ngược $0\le T\le t$, $|y-x|\le a(t-T)$; đổi thứ tự tích phân trên miền này cho dạng tích phân Green tương đương.

\textbf{Bước 4. Dạng Green trực tiếp.} Trên tam giác đó, đặt $p=u_T$, $q=-a^2u_y$, nên $p_T+q_y=F$. Tích phân và lấy thông lượng hai cạnh đặc trưng, trên đó $du/dT=u_T\mp au_y$, cho
\[\int_0^t\int_{x-a(t-T)}^{x+a(t-T)}F\,dy\,dT
=2au(x,t)-a[f(x-at)+f(x+at)]-\int_{x-at}^{x+at}g(y)dy.\]
Chia $2a$ thu đúng công thức. Công thức này đồng thời chỉ ra dấu và hệ số của từng dữ liệu.

\textbf{Chi tiết thông lượng trên từng cạnh tam giác Green.} Đặt $\ell(T)=x-a(t-T)$, $r(T)=x+a(t-T)$ và $M(T)=\int_{\ell(T)}^{r(T)}u_T(y,T)dy$. Quy tắc Leibniz cho
\[M'=\int_\ell^r u_{TT}dy-a u_T(r,T)-a u_T(\ell,T).\]
Tích phân $u_{TT}-a^2u_{yy}=F$ theo $y$, rồi thay dòng trên:
\[\int_\ell^r Fdy=M'+a(u_T+a u_y)(\ell,T)+a(u_T-a u_y)(r,T).\]
Cạnh trái có $\ell'=a$, nên $(u_T+a u_y)(\ell,T)=d[u(\ell(T),T)]/dT$; cạnh phải có $r'=-a$, nên $(u_T-a u_y)(r,T)=d[u(r(T),T)]/dT$. Tích phân mỗi cạnh từ $0$ đến $t$ cho
\[a[u(x,t)-f(x-at)],\qquad a[u(x,t)-f(x+at)].\]
Mặt đáy cho $M(t)-M(0)=0-\int_{x-at}^{x+at}g(y)dy$, vì ở đỉnh đoạn tích phân co về độ dài $0$. Cộng ba đóng góp được đẳng thức Green ở bước 4, rồi chia $2a$ thu half-sum, tích phân vận tốc và tích phân nguồn.

%ID BT_3_1_3_05
\textbf{Bước 1. Giải phần thuần nhất.} Thay $f,g$ vào d'Alembert. Nguyên hàm cần dùng lần lượt là $-\cos2s/2$, $\arctan s$, $-\cos s$, $e^s$, $\sin s$ và $H$ của xung tam giác. Sau đó tìm phần nguồn có hai dữ liệu đầu bằng không.

\textbf{Bước 2. Tính từng nguồn.}

(i) $a=2$, nguồn $x\cos t$ có nghiệm riêng $x(1-\cos t)$ vì đạo hàm không gian cấp hai bằng không. Suy
\[u=e^{-x}\cosh2t+\tfrac14\sin2x\sin4t+x(1-\cos t).\]
(ii) $a=3$. Nguồn $e^x$ dùng $e^xq$, $q''-9q=1$, $q(0)=q'(0)=0$, cho $q=(\cosh3t-1)/9$. Nguồn $\sin t$ độc lập $x$ cho $t-\sin t$:
\[u=\frac{\arctan(x+3t)-\arctan(x-3t)}6+\frac{e^x(\cosh3t-1)}9+t-\sin t.\]
(iii) $a=4$, nguồn $x+t$ cho $xt^2/2+t^3/6$:
\[u=x^2+16t^2+\tfrac14\sin x\sin4t+xt^2/2+t^3/6.\]
(iv) $a=7$, phần vận tốc $e^x$ là $e^x\sinh7t/7$, phần nguồn $xt$ là $xt^3/6$.

(v) $a=1$, $F=\sin(x+t)=\sin x\cos t+\cos x\sin t$. Hai ODE là $q''+q=\cos t$ và $r''+r=\sin t$ với dữ liệu không, ở cộng hưởng. Suy $q=t\sin t/2$, $r=(\sin t-t\cos t)/2$:
\[u=x+\cos x\sin t+\tfrac12\{\sin x\,t\sin t+\cos x(\sin t-t\cos t)\}.\]
(vi) $a=2$, $f$ là hình thang trên $[0,3]$ với nhánh $s,1,3-s$ tại các đoạn $[0,1],[1,2],[2,3]$. $H$ là nguyên hàm tam giác ở Bài 2(vi). Nguồn $-e^x$ cho
\[u=\tfrac12[f(x-2t)+f(x+2t)]+\frac{H(x+2t)-H(x-2t)}4-\frac{e^x(\cosh2t-1)}4.\]
\textbf{Bước 3. Kiểm tra.} Mỗi phần nguồn và đạo hàm thời gian của nó bằng không ở $t=0$. Các ODE vừa lập chứng minh PDE, không chỉ đoán dạng đáp án. Câu (vi) có góc của $f,g$, nên nghiệm được hiểu từng miền/ theo nghĩa yếu trên các đường $x\pm2t=0,1,2,3$ và $x\pm2t=-1,0,1$.

\textbf{Bước 4. Vẽ đặc trưng và các lát cắt.} Với mỗi tốc độ $a$ của từng câu, các đường xuất phát từ $(b,0)$ là $x=b+at$ và $x=b-at$. Nguồn thay đổi giá trị trên miền ảnh hưởng nhưng không đổi hai họ đặc trưng này. Hình kèm lấy $t=0,0.15,0.35,0.65$ và thay trực tiếp vào sáu công thức đã tính; hình thứ hai cho mặt phẳng $x,t$ với các đường đặc trưng mẫu.

\textbf{Chi tiết giải các ODE nguồn.} (i) Đặt $u_F=xq$, có $q''=\cos t$, $q'(0)=q(0)=0$. Tích phân $q'=\sin t$, $q=1-\cos t$. (ii) Với $u_F=e^xq$, $q''-9q=1$: thử $q_p=-1/9$, phần thuần nhất $A\cosh3t+B\sinh3t$. Dữ liệu cho $A=1/9,B=0$. Nguồn $\sin t$ cho $r''=\sin t$, $r'=1-\cos t$, $r=t-\sin t$.

(iii) $q''=x+t$, tích phân theo thời gian với $x$ cố định được $q'=xt+t^2/2$, $q=xt^2/2+t^3/6$. (iv) $q''=xt$ cho $q'=xt^2/2$, $q=xt^3/6$.

(v) Tính trực tiếp $q=t\sin t/2$: $q'=(\sin t+t\cos t)/2$, $q''=\cos t-t\sin t/2$, nên $q''+q=\cos t$. Với $r=(\sin t-t\cos t)/2$, $r'=t\sin t/2$, $r''=(\sin t+t\cos t)/2$, nên $r''+r=\sin t$. Cả $q,r$ và đạo hàm đầu bằng $0$. (vi) $u_F=e^xq$ cho $q''-4q=-1$, $q_p=1/4$, $q=A\cosh2t+B\sinh2t+1/4$; dữ liệu cho $A=-1/4,B=0$.

%ID BT_3_1_3_06
\textbf{Bước 1. Đạo hàm công thức nửa trục.} Trong miền $x<at$, đặt $s=x+at$, $r=at-x$, $\tau=t-x/a$. Từ $u=h(\tau)+(f(s)-f(r))/2+(2a)^{-1}\int_r^s g$ thu
\[u_x=-h'(\tau)/a+\tfrac12(f'(s)+f'(r))+\frac{g(s)+g(r)}{2a}.\]
Đây là công thức (3.1.13); dấu cộng của $g(r)$ do cận dưới có đạo hàm $r_x=-1$.

\textbf{Bước 2. Đạo hàm thời gian hai lần.}
\[u_{tt}=h''(\tau)+\tfrac{a^2}2(f''(s)-f''(r))+\tfrac a2(g'(s)-g'(r)).\]
Ngoài mặt sóng, thay dấu trừ giữa hai $f''$ thành dấu cộng, và giữ cùng hiệu $g'$. Tại $x=at$, hiệu hai vết là $h''(0)-a^2f''(0)$, cho công thức ghép (3.1.15).

\textbf{Bước 3. Kết luận.} Cần $h(0)=f(0)$ để liên tục, $h'(0)=g(0)$ để ghép cấp một, và $h''(0)=a^2f''(0)$ để ghép cấp hai. Những đạo hàm tính ở đây cũng cho $u_{tt}=a^2u_{xx}$ trong miền phản xạ và đúng biên $u(0,t)=h(t)$.

\textbf{Chi tiết lần đạo hàm thứ nhất theo thời gian.} Với $s_t=r_t=a,\tau_t=1$,
\[u_t=h'(\tau)+\frac a2[f'(s)-f'(r)]+\frac12[g(s)-g(r)].\]
Lấy tiếp đạo hàm: $h'$ cho $h''$, hai hạng $f'$ cho $a^2(f''(s)-f''(r))/2$, hai hạng $g$ cho $a(g'(s)-g'(r))/2$. Tương tự, lấy đạo hàm $x$ của công thức $u_x$ cho
\[u_{xx}=h''(\tau)/a^2+(f''(s)-f''(r))/2+(g'(s)-g'(r))/(2a),\]
chứng minh trực tiếp $u_{tt}=a^2u_{xx}$.

%ID BT_3_1_3_07
\textbf{Bước 1. Mở rộng lẻ.} Đặt $\widetilde f(-s)=-f(s)$, $\widetilde g(-s)=-g(s)$. So vết tại $0$: $f(0)=0$ nối giá trị, $f'$ tự nối, $f''(0)=0$ nối cấp hai; $g(0)=0$ nối giá trị, $g'$ tự nối. Do đó $\widetilde f\in C^2$, $\widetilde g\in C^1$ dưới các điều kiện tương thích.

\textbf{Bước 2. Giải trên toàn trục.} Đặt $H(s)=\int_0^s\widetilde g(r)dr$; $H$ chẵn. D'Alembert cho $u=(\widetilde f(x-at)+\widetilde f(x+at))/2+(H(x+at)-H(x-at))/(2a)$. Tại $x=0$, hai giá trị $\widetilde f$ đối nhau và hai giá trị $H$ bằng nhau, nên $u(0,t)=0$.

\textbf{Bước 3. Viết lại bằng dữ liệu dương.} Khi $x\ge at$ giữ nguyên công thức. Khi $0<x<at$, $\widetilde f(x-at)=-f(at-x)$, còn $H(at+x)-H(at-x)=\int_{at-x}^{at+x}g$. Thu công thức phản xạ trong phần kiến thức với $h=0$. Ranh giới đúng là $x=at$; $x=t$ chỉ đúng nếu $a=1$. Dữ liệu đầu và tính $C^2$ suy trực tiếp từ mở rộng.

%ID BT_3_1_3_08
\textbf{Bước 1. Dùng công thức Dirichlet với thời điểm trễ.} Ở $x\ge at$ sóng chưa chạm biên; ở $x<at$ dùng $h(t-x/a)$ và phản xạ lẻ. Dưới đây giữ hai miền khi biên làm thay đổi vận tốc.

(i) $a=2$, $f=\sin x$, $g=\cos2x$, $h=0$. Mở rộng $g$ lẻ, $H(s)=\sin(2|s|)/2$. Do đó
\[u=\sin x\cos2t+\begin{cases}\cos2x\sin4t/4&x\ge2t,\\\sin2x\cos4t/4&0<x<2t.\end{cases}\]
$g(0)=1\ne h'(0)$ nên không $C^1$ trên mặt sóng. Sau khi mặt sóng đi qua mỗi $x$, công thức tuần hoàn chu kỳ $\pi$; toàn nghiệm không có chu kỳ chung vì hai nhánh khác nhau.

(ii) $a=3$, $f=g=0$, $h=t/(1+t)$. Đặt $\tau=t-x/3$: $u=0$ khi $\tau\le0$, $u=\tau/(1+\tau)$ khi $\tau>0$. Trên tia $x=ct$, giới hạn là $1$ nếu $0\le c<3$, và $0$ nếu $c\ge3$. Nghiệm không tuần hoàn: tại mỗi $x$ cố định, sau khi mặt sóng đến, $u$ tăng nghiêm ngặt về $1$.

(iii) $a=1$, mở rộng lẻ $f$ là $s|s|$, còn $g=\sin s$ đã lẻ:
\[u=\tfrac12[(x+t)|x+t|+(x-t)|x-t|]+\sin x\sin t.\]
Nghiệm $C^1$, không $C^2$ trên $x=t$ vì $f''(0)=2\ne0$, và không tuần hoàn.

(iv) $a=5$, $f=0$, $g=\sin x-\sin3x/3$ có $H=-\cos s+\cos3s/9$. Suy $u=\sin x\sin5t/5-\sin3x\sin15t/45$, trơn, chu kỳ $2\pi/5$.

(v) $a=2$, $f=\cos x-\sin x$, $g=0$, $h=\cos2t-\sin2t$. Ngoài mặt sóng, half-sum của $f$ là $(\cos x-\sin x)\cos2t$. Trong mặt sóng:
\[h(t-x/2)+\tfrac12[f(x+2t)-f(2t-x)]=\cos x(\cos2t-\sin2t).\]
Vậy
\[u=\begin{cases}(\cos x-\sin x)\cos2t&x\ge2t,\\\cos x(\cos2t-\sin2t)&0<x<2t.\end{cases}\]
Công thức cũ $\cos(x-2t)-\sin(x+2t)$ không thỏa $u_t(x,0)=0$. Nghiệm đúng liên tục, không $C^1$ vì $h'(0)=-2\ne0$. Chỉ nhánh sau mặt sóng có chu kỳ $\pi$ ở từng điểm.

(vi) $a=1$, $g=\chi_{[1,2]}-\chi_{[3,4]}$. Mở rộng lẻ có nguyên hàm chẵn $H(s)=[|s|-1]_0^1-[|s|-3]_0^1$, nên $u=(H(x+t)-H(x-t))/2$. Tại $x$ cố định, hai giá trị $H$ đều bằng không sau $t>x+4$, nên $u=0$; đây là xung quá độ, không tuần hoàn. Nghiệm liên tục nhưng có gãy cấp một khi các cận gặp $\pm1,\pm2,\pm3,\pm4$.

\textbf{Bước 2. Kiểm tra chung.} Thế $x=0$ vào nhánh trong cho đúng $h$; thế $t=0$ vào nhánh ngoài cho $f,g$. So hai vết trên $x=at$ theo ba điều kiện tương thích, không dùng tính tuần hoàn của nhánh trong để kết luận tuần hoàn toàn miền.

\textbf{Chi tiết hai nhánh vận tốc ở (i).} Ngoài mặt sóng,
\[\frac14\int_{x-2t}^{x+2t}\cos2s\,ds
=\frac{\sin(2x+4t)-\sin(2x-4t)}8=\frac14\cos2x\sin4t.\]
Trong mặt sóng, cận dưới là $2t-x$, nên hiệu sine bằng $\sin(4t+2x)-\sin(4t-2x)=2\cos4t\sin2x$; chia $8$ cho nhánh đã ghi.

\textbf{Chi tiết rút gọn (v).} Trong miền phản xạ, biên là $\cos(2t-x)-\sin(2t-x)$. Half-difference của $f$ bằng
\[-\sin2t\sin x-\cos2t\sin x.\]
Khai triển biên thành $\cos2t\cos x+\sin2t\sin x-\sin2t\cos x+\cos2t\sin x$ rồi cộng, hai hạng chứa $\sin x$ triệt tiêu, còn $\cos x(\cos2t-\sin2t)$.

%ID BT_3_1_3_09
\textbf{Bước 1. Tìm các hàm sóng ở hai miền.} Biên nghiêng là $t=\alpha x$, không phải $x=0$. Viết $u=p+A(x+at)+B(x-at)$, trong đó $p$ là nghiệm riêng. Dữ liệu ở $t=0$ xác định $A,B$ trên đối số dương; dữ liệu biên xác định $B$ ở đối số âm.

(i) $a=1$, $\alpha=2$, $F=0$, $f=x$, $g=0$, vết biên $3x$. Trên $s>0$, $A(s)=B(s)=s/2$ (bỏ hai hằng số triệt tiêu). Tại $t=2x$, $A(3x)+B(-x)=3x$, nên $B(-x)=3x/2$, tức $B(s)=-3s/2$ khi $s<0$. Suy
\[u=\begin{cases}x&0\le t\le x,\\2t-x&x<t<2x.\end{cases}\]
Hai vết giá trị bằng nhau, nhưng $u_t$ đổi từ $0$ sang $2$; không có nghiệm $C^1$ trên toàn tam giác thỏa mọi dữ liệu này.

(ii) $a=2$, $\alpha=3$, $F=1$, $f=g=0$, biên bằng không. Chọn $p=-x^2/8$, vì $p_{tt}-4p_{xx}=1$. Dữ liệu đầu cho $A=B=s^2/16$ trên $s>0$. Tại $t=3x$, $A(7x)+B(-5x)-x^2/8=0$, suy $B(s)=-47s^2/400$ khi $s<0$:
\[u=\begin{cases}t^2/2&0\le t\le x/2,\\-x^2/8+(x+2t)^2/16-47(x-2t)^2/400&x/2<t<3x.\end{cases}\]
\textbf{Bước 2. Kiểm tra.} Tại đường $x=2t$, $B$ và $B'$ nối bằng không nên nghiệm $C^1$, nhưng $B''$ đổi từ $1/8$ sang $-47/200$, nên không $C^2$. Mỗi nhánh thỏa PDE; tại $t=3x$ đa thức triệt tiêu đúng. Hai kết quả phải hiểu từng miền khi đề yêu cầu mức trơn cao hơn mức dữ liệu cho phép.

%ID BT_3_1_3_10
\textbf{Bước 1. Điều kiện Neumann.} $u_x(0,t)=0$ kéo theo $f'(0)=0$, $g'(0)=0$ nếu nghiệm $C^2$. Mở rộng chẵn của $f,g$ có vết $f'$ và $g'$ đổi dấu, nên hai điều kiện này bảo đảm các vết cần thiết nối nhau.

\textbf{Bước 2. Áp dụng d'Alembert.} Với mở rộng chẵn $\widetilde f,\widetilde g$, nguyên hàm $H=\int_0^s\widetilde g$ là lẻ. Nghiệm là half-sum cộng hiệu $H/(2a)$. Nó chẵn theo $x$, do đó đạo hàm không gian ở $0$ bằng không. Trong miền $x<at$ có thể viết trực tiếp
\[u=\frac{f(at-x)+f(at+x)}2+\frac1a\int_0^{at-x}g(s)ds+\frac1{2a}\int_{at-x}^{at+x}g(s)ds.\]
Hạng tích phân đầu xuất hiện vì đoạn $[x-at,0]$ cũng góp vận tốc sau phản xạ chẵn.

\textbf{Bước 3. Ví dụ.} Với $a=2$, $f=x^2$, $g=\cos x$, cả hai đã chẵn và tương thích. Half-sum cho $x^2+4t^2$, nguyên hàm $\sin s$ cho $\cos x\sin2t/2$. Vậy $u=x^2+4t^2+\cos x\sin2t/2$. Kiểm tra $u_x(0,t)=0$, $u_t(x,0)=\cos x$ và $u_{tt}=4u_{xx}$ trực tiếp.

%ID BT_3_1_3_11
\textbf{Bước 1. Chuyển đạo hàm thành bài Dirichlet.} Đặt $v=u_x$. Khi dữ liệu đủ trơn, $v_{tt}=a^2v_{xx}$, $v(0,t)=h(t)$, $v(x,0)=f'(x)$, $v_t(x,0)=g'(x)$. Tương thích cấp hai cho bài này là $h(0)=f'(0)$, $h'(0)=g'(0)$, $h''(0)=a^2f'''(0)$.

\textbf{Bước 2. Tích phân lại và giữ hằng số theo thời gian.} $u(x,t)=C(t)+\int_0^xv(s,t)ds$. Thế vào PDE cho $C''(t)=a^2v_x(0,t)$, $C(0)=f(0)$, $C'(0)=g(0)$. Không được bỏ $C$ chỉ vì đã biết $u_x$.

\textbf{Bước 3. Công thức tiện dụng tương đương.} Gọi $u_e$ là nghiệm d'Alembert của các mở rộng chẵn $f,g$. Đặt
\[u=u_e-a\int_0^{(t-x/a)_+}h(s)ds.\]
Ở $x<at$, đạo hàm không gian của hạng cuối là $h(t-x/a)$, cho đúng $u_x(0,t)=h(t)$; hạng đó là sóng chạy nên thỏa PDE. Nếu $h(0)\ne0$, góc của $u_e$ do $f'(0)$ được hạng biên bù khi tương thích đúng.

\textbf{Bước 4. Viết hai nhánh.} Ngoài mặt sóng dùng d'Alembert với dữ liệu dương. Trong mặt sóng dùng công thức Bài 10 rồi trừ $a\int_0^{t-x/a}h$. Hai dữ liệu đầu giữ nguyên vì hạng bổ sung chưa đến $x>0$ tại $t=0$. Nếu chỉ cần $u\in C^2$, hai tương thích đầu đủ với mức trơn phù hợp; điều kiện thứ ba ở trên phục vụ $v\in C^2$, tức mức trơn cao hơn của $u$.

%ID BT_3_1_3_12
\textbf{Bước 1. Tách hai sóng trên dữ liệu gốc.} Đặt $H(s)=\int_0^sg$, $A=f/2-H/(2a)$, $B=f/2+H/(2a)$ trên $0\le s\le L$. Hai biên là $u(0,t)=\alpha(t)$ và $u(L,t)=\beta(t)$. Chỉ xét $t\le L/a$, nên mỗi chân đặc trưng phản xạ tối đa một lần.

\textbf{Bước 2. Viết đủ bốn miền.} Đặt $\xi=x-at$, $\eta=x+at$:
\[u=\begin{cases}
A(\xi)+B(\eta)&\xi\ge0,\ \eta\le L,\\
\alpha(t-x/a)-B(at-x)+B(x+at)&\xi<0,\ \eta\le L,\\
A(x-at)+\beta(t-(L-x)/a)-A(2L-x-at)&\xi\ge0,\ \eta>L,\\
\alpha(t-x/a)+\beta(t-(L-x)/a)-B(at-x)-A(2L-x-at)&\xi<0,\ \eta>L.
\end{cases}\]
Hai công thức phản xạ suy từ $A(-s)=\alpha(s/a)-B(s)$ và $B(L+s)=\beta(s/a)-A(L-s)$; mọi đối số của $A,B$ còn lại đều trong $[0,L]$.

\textbf{Bước 3. Trường hợp (i).} $a=2$, $L=\pi$, $f=\sin x$, $g=1$, $\alpha=\beta=0$. $A=\sin s/2-s/4$, $B=\sin s/2+s/4$. Trong bốn miền theo đúng thứ tự trên:
\[u=\sin x\cos2t+\begin{cases}t&\text{miền I},\\x/2&\text{miền II},\\(\pi-x)/2&\text{miền III},\\\pi/2-t&\text{miền IV}.\end{cases}\]
Phần vận tốc liên tục, nhưng có góc trên hai mặt sóng vì $g$ không tương thích với biên không.

\textbf{Bước 4. Trường hợp (ii).} $a=L=1$, $f=g=0$, $\alpha(t)=\sin\pi t$, $\beta(t)=t$. Vì $A=B=0$, bốn nhánh là $0$, $\sin\pi(t-x)$, $t-1+x$, và $\sin\pi(t-x)+t-1+x$. Thế $x=0,1$ trong khoảng $0\le t\le1$ cho đúng hai biên. Vận tốc góc không tương thích ($\alpha'(0)=\pi$, $\beta'(0)=1$), nên không có nghiệm $C^1$ trên toàn miền; các nhánh và đường ghép đã mô tả chính xác sự mất trơn.

\textbf{Chi tiết thế vào miền cả hai mặt sóng đã phản xạ ở (i).} Đặt $p=2t-x$, $q=2\pi-x-2t$. Khi $\xi<0,\eta>\pi$, $p,q\in[0,\pi]$ và $u=-B(p)-A(q)$:
\[u=-\frac{\sin(2t-x)}2-\frac{\sin(2\pi-x-2t)}2+\frac{q-p}{4}
=\frac{\sin(x+2t)+\sin(x-2t)}2+\frac{2\pi-4t}{4}.\]
Tổng sine bằng $2\sin x\cos2t$, hạng tuyến tính bằng $\pi/2-t$. Ở miền chỉ phản xạ trái, $-B(2t-x)+B(x+2t)$ cho hạng tuyến tính $[(x+2t)-(2t-x)]/4=x/2$. Ở miền chỉ phản xạ phải, $A(x-2t)-A(2\pi-x-2t)$ cho $[-(x-2t)+(2\pi-x-2t)]/4=(\pi-x)/2$.
